\subsection{Structural geometric unibranchness under actual local ring equivalences}

\begin{theorem}[The original reduction and normalization equivalences]
\label{algebra:actual-unibranch-normalization-ring-equivalence}
\lean{EnvelopeProperty.LocalNormalization.map_nilradical_eq,
EnvelopeProperty.LocalNormalization.reducedEquiv,
EnvelopeProperty.LocalNormalization.reducedEquiv_mk,
EnvelopeProperty.LocalNormalization.fractionEquiv,
EnvelopeProperty.LocalNormalization.fractionEquiv_algebraMap,
EnvelopeProperty.LocalNormalization.normalizationEquiv,
EnvelopeProperty.LocalNormalization.normalizationEquiv_coe,
EnvelopeProperty.LocalNormalization.normalizationEquiv_coefficientMap}
\uses{algebra:actual-structural-unibranch-local-ring}
An ACTUAL ring equivalence $e:R\simeq S$ induces genuine whole ring
equivalences between their ORIGINAL reductions, their actual reduced
fraction rings and their LITERAL reduced normalizations $N(R)\simeq N(S)$.
The entire original coefficient square commutes:
$e_N\circ c_R=c_S\circ e$ on EVERY original ring element.
No domain or localness is needed for these algebraic equivalences.
\leanok
\end{theorem}
\begin{proof}
True nilpotent preservation under $e$ and $e^{-1}$ proves equality of
the mapped ORIGINAL nilradical with the original target nilradical.
The actual quotient universal property produces the genuine reduction
equivalence $e_{\mathrm{red}}$, retaining each original residue class.
Actual fraction-ring uniqueness gives its whole ambient fraction-ring
equivalence, preserving the true reduced coefficient maps.

Map a genuine monic integral relation through the ACTUAL reduction
and ambient fraction-ring homomorphisms. Their whole coefficient square
shows that the mapped monic polynomial annihilates the entire mapped
element. Thus every original integral element remains integral over
the target reduction. Applying the SAME argument to the inverse
ring equivalences gives the genuine reverse map. Restrict the whole
ambient equivalence to these actual integral elements to obtain
$e_N$. Actual quotient/fraction coefficient identities prove the
ENTIRE original coefficient square after embedding into the actual
target ambient fraction ring. Both inverse identities come from the
true whole ambient inverse laws; no unrelated scalar action,
normalization or localization certificate is supplied.
\end{proof}

\begin{theorem}[Standard structural local properties really transport]
\label{algebra:actual-structural-unibranch-ring-equivalence}
\lean{EnvelopeProperty.IsUnibranchLocalRing.of_ringEquiv,
EnvelopeProperty.LocalNormalization.residueFieldMap_ringEquiv_square,
EnvelopeProperty.IsGeometricallyUnibranchLocalRing.of_ringEquiv}
\uses{algebra:actual-unibranch-normalization-ring-equivalence,algebra:actual-normalization-local-residue-square}
For actual local rings $R,S$ and a genuine equivalence $e:R\simeq S$,
standard unibranchness of $R$ implies standard unibranchness of $S$.
The ACTUAL residue maps fit the whole commuting square with the
residue equivalences induced by $e$ and $e_N$. Thus standard
GEOMETRIC unibranchness also transports to $S$.
\leanok
\end{theorem}
\begin{proof}
The true reduction equivalence transports the actual domain property,
and $e_N$ transports actual localness of the literal normalization.
This constructs the standard target structural class rather than
assuming its spectrum fibre property. Both original coefficient maps
are consequently genuine local maps, derived in the prior structural
core. The whole coefficient square above and the ACTUAL residue
functor composition law produce equality on entire residue fields;
quotient surjectivity proves it from EVERY original representative.

Solve the entire residue square for the target coefficient residue
map using the two genuine residue equivalences. Each equivalence
(and its inverse) is surjective, hence purely inseparable. Compose
these actual field maps with the source standard pure-inseparable
residue map. The resulting map is exactly the target's actual
coefficient residue map, establishing its standard geometric class.
There is no target radicial, native-map or singleton-fibre hypothesis.
\end{proof}

\begin{theorem}[Genuine scheme stalks supply the actual affine normalization theorem]
\label{scheme:actual-structural-unibranch-stalk-normalization}
\lean{EnvelopeProperty.IsGeometricallyUnibranchScheme,
EnvelopeProperty.isGeometricallyUnibranch_atPrime_of_scheme,
EnvelopeProperty.DomainNormalization.projection_universallyInjective_of_schemeStalks,
EnvelopeProperty.DomainNormalization.nativeHomeomorph_of_schemeStalks,
EnvelopeProperty.DomainNormalization.nativeHomeomorph_of_schemeStalks_apply,
EnvelopeProperty.DomainNormalization.relativeNativeHomeomorph_of_schemeStalks,
EnvelopeProperty.DomainNormalization.relativeNativeHomeomorph_of_schemeStalks_apply}
\uses{algebra:actual-structural-unibranch-ring-equivalence,scheme:structural-domain-normalization-native-homeomorphism}
The standard scheme property is geometric unibranchness of EVERY
ACTUAL local stalk. For an actual integral affine scheme
$\operatorname{Spec}A$ with this property, its literal original
normalization has the proved native homeomorphism, absolutely and
above ANY exact common prescribed coefficient seminorm. The
homeomorphism retains the WHOLE original normalization point map.
\leanok
\end{theorem}
\begin{proof}
For every original $p\in\operatorname{Spec}A$, Mathlib's TRUE
$\operatorname{Spec}A$ stalk isomorphism identifies its ACTUAL
scheme stalk with the ORIGINAL localization $A_p$. Apply the
preceding standard structural invariance theorem to this GENUINE
ring equivalence. It derives standard geometric unibranchness of
that original $A_p$, including the actual residue-field map.
The already proved integral-affine normalization theorem therefore
applies with these DERIVED local properties. It constructs the
literal normalization's actual scheme radiciality and the native
homeomorphisms retaining their ENTIRE original maps. No local-ring,
residue comparison or native-bijection certificate is requested.
General nonintegral scheme normalization gluing and the distinct
strict-henselization characterization remain separate work.
\end{proof}
